% Source: https://artofproblemsolving.com/community/c3h1766239_some_imaginary_math
% Stable identity: aops-topic-1766239-q01
The absolute value of $a+bi$ is its distance from $0$ in the complex plane, $|a+bi|=\sqrt{a^2+b^2}$. So
\[
|5-12i|=\sqrt{5^2+(-12)^2}=\sqrt{169}=\boxed{13}.
\]
The student's steps are wrong in two ways. First, $12\sqrt{-1}=12i$ is not a real number, and $-12i$ cannot be rewritten as $+\sqrt{12}$ (that step treats $\sqrt{-1}$ as $-1$ and also replaces $12$ by $\sqrt{12}$). Second, $5-12i$ itself is not a real number at all, so its absolute value is not found by simplifying it to a real number; it is the length $\sqrt{a^2+b^2}$.
